\section{Ordinary and $\Lambda$-adic passage}\label{sec:stage-04}
% Stage 04 (lambda-adic). Drafter: Opus stage drafter, 2026-11-14. Source: seed \S7, \S2.6, \S3.3.
% Local macros (not in preamble.tex; requested in missing/stage-04.md):
\providecommand{\NUp}{N_{U_p}}                 % the U_p normalisation fixed in stage 01
\providecommand{\Kp}[1]{K_{#1}}                % K_s = U_1(p^s)_p
\providecommand{\Kpp}[1]{K'_{#1}}              % K'_s = U_1(p^s)_p \cap U_0(p^{s+1})_p
\providecommand{\TordB}{\T^{\mathrm{ord}}_B}   % ordinary definite Hecke algebra (NOT \Tord, which is Hida's Tate module)

Seed sources: seed \S7 (Lemma ekills, Prop.~free, Thm.~lambdaribet, the pairing paragraph), \S2.6 (Hida
notation), \S3.3 (nearly ordinary remark). Conventions are those of stage~01, in particular:
$U_p=[U\varpi_pU]$ with $\varpi_p=\iota_p^{-1}\mathrm{diag}(p,1)$, $U_p[x]=\sum_{j=0}^{p-1}[x\varpi_j]$,
$\varpi_j=\iota_p^{-1}\left(\begin{smallmatrix}p&j\\0&1\end{smallmatrix}\right)$ (normalisation $\NUp$, \emph{not} the seed's literal
$\mathrm{diag}(1,p)$); the adelic diamond $\langle d\rangle$ (right translation by $\delta_d$, lower right entry $d$) is the
\emph{inverse} of the classical one; the norm map $\nu$ intertwines operators with their adjoints; freeness in Lemma stab is for
$\Delta_s$ only; the weighted pairing on $\Vk$ is perfect over $\Zp$ only for $p>k-2$.

%----------------------------------------------------------------------
\subsection{Goals}
%----------------------------------------------------------------------
\begin{enumerate}[label=(G\arabic*)]
\item \emph{Lemma ekills.} Prove that the ordinary projector $e$ kills the degree term, the vertex term of the DR graph, the
  combinatorial component group and $H_0(\GU,\Zp)$, and that the Brandt pairing restricts to a perfect pairing on ordinary
  parts. Make precise \emph{which} $e$ appears on each side of the pairing (pushforward-compatible $e$ versus its adjoint).
\item \emph{Hida theory on the definite side.} Define $\Mord(U)=\varprojlim_se\,\Zp[\XB{UU_s}]$, prove its $\Lambda$-freeness and
  control by an \emph{elementary} two-step argument (Hida's lemma plus a free $\Delta$-cover), and record the cited statements
  ([Hs, Thm.~4.2], [LV11], Hida 1988). Fix the $\Lambda$-module structure and the weight-$k$ arithmetic ideals in the
  conventions of stage~01 (the seed's $P_k$ is off by the diamond inversion).
\item \emph{$\Lambda$-adic definite Ribet sequence.} Prove Thm.~lambdaribet with the $U_p$ normalisation $\NUp$, the rank formula,
  and $\Lambda$-freeness of $\Xord_{Dq}(U)$.
\item \emph{Pairing, Gorenstein property and $\eta^B_F$.} State Hsieh's twisted $\Lambda$-adic pairing, the freeness of
  $\Mord{}_{,\m}$ over the Gorenstein ring $\TordB{}_{,\m}$ under an \emph{explicit} residual hypothesis, and $\eta^B_F$ as a
  self-pairing, all translated into the seed's conventions.
\item \emph{Nearly ordinary remark.} Make seed \S3.3 precise on the definite side.
\end{enumerate}

%----------------------------------------------------------------------
\subsection{Definitions and statements}
%----------------------------------------------------------------------

\subsubsection{Standing notation}\label{ss:s04-notation}
$p\ge5$, $p\nmid D$; $\cO=\cOE$ for a finite extension $E/\Qp$, uniformiser $\varpi$. $\Lam=\cOE[[1+p\Zp]]$, $[u]$ the group
element of $u\in1+p\Zp$. A \emph{level prime to $p$} is a compact open $U=U^{(p)}\times\GL_2(\Zp)\subseteq\hat\cO^\times$ with
$\Nrd U=\Zhat^\times$. For such $U$ and $s\ge0$ we write
\[UU_s:=U^{(p)}\times\Kp{s},\qquad \Kp{s}:=\iota_p^{-1}\Uone{p^s}_p,\qquad \Kpp{s}:=\Kp{s}\cap\iota_p^{-1}\Uzero{p^{s+1}}_p ,\]
so $UU_s=U_s=\Uone{Mp^s}$ when $U=\Uone{M}$ (this fixes the seed's ``$UU_s$''). Put $X_s(U):=\XB{UU_s}$, $X'_s(U):=\XB{U^{(p)}\Kpp{s}}$.
For $s\ge1$ every $\Gamma_x$ at these levels is trivial and $-1\notin UU_s$ (\texttt{if:s01-classset}, Lemma stab (b)), so
$w_x=1$ and $\nu=\mathrm{id}$ on $\Zp[X_s]$. $\Delta_s=(1+p\Zp)/(1+p^s\Zp)$ acts on $X_s(U)$ by adelic diamonds $[x]\mapsto[x\delta_d]$,
freely (Lemma stab (b)); the full group $(\Z/p^s)^\times=\mu_{p-1}\times\Delta_s$ acts, but not necessarily freely.

All modules below are homological ($\Zp[X]$, pushforward) unless marked otherwise, with $U_p[x]=\sum_j[x\varpi_j]$ and Hecke
operators acting as transposes of the cohomological ones (\texttt{if:s01-brandt}). $\pi_{s+1,s}:X_{s+1}\to X_s$ is the natural
projection and $\pi_*$ the pushforward.

\begin{definition}[ordinary projector; seed \S2.6]\label{def:s04-e}
For a finitely generated $\Zp$-module $M$ with an endomorphism $T$, $e_T:=\lim_{n\to\infty}T^{n!}$. It exists, is an idempotent
in the closed subalgebra $\Zp[[T]]\subseteq\End M$ generated by $T$, and $M=e_TM\oplus(1-e_T)M$ with $T$ invertible on $e_TM$ and
topologically nilpotent on $(1-e_T)M$. The functor $M\mapsto e_TM$ is exact on $\Zp[T]$-modules of this kind. If $M$ is finite,
$e_T=T^{N}$ for any $N$ that is a multiple of the period and at least the index of $T$ in the finite monoid $\End M$
(e.g.\ $N=|\End M|!$). We write $e:=e_{U_p}$ and $e^*:=e_{U_p^*}$, where $U_p^*=\langle p\rangle_MU'_p$ is the
$\langle,\rangle$-adjoint (\texttt{if:s01-hecke}).
\end{definition}
\noindent\stknown{} \lit{hida\_1993\_elementary-theory-of-l-functions-and-eisenstein-series}{\S7.2, Lemma~1, p.~201: for a commutative
$\cO$-algebra $A$ of finite rank and $x\in A$, $\lim x^{n!}$ exists and is idempotent} (apply to $A=\cO[T]\subseteq\End M$);
the finite-monoid statement is elementary \sttoprove.

\subsubsection{Lemma ekills}

\begin{lemma}[$e$ kills degree, vertex and component terms; seed Lemma ekills]\label{lem:s04-ekills}
Let $s\ge1$, $V=\Zp$, $U$ of level prime to $p$.
\begin{enumerate}[label=(\alph*)]
\item $\deg:\Zp[X_s(U)]\to\Zp$ satisfies $\deg\circ U_p=p\cdot\deg$. Hence $e$ kills $\Zp$ (with $U_p=p$), $e\,\Zp[X_s]^0=e\,\Zp[X_s]$,
  and, for $q\nmid Dp$ and $U_q=\GL_2(\Zq)$, $e$ kills $H_0(\Gamma_{UU_s},\Zp)\cong\Zp$ (with the $U_p$-action induced from
  $\coker\partial$, \texttt{if:s02-coh}).
\item ($B=B_q$, DR graph $\GDR_q(UU_s)$, \texttt{if:s03-graph}.) Let $U_p$ act on the vertex module $\Zp[\{v_0,v_1\}]$ by $p$; then
  $\partial$ is $U_p$-equivariant, $e$ kills the vertex term, and $e\,\Xcomb_q(UU_s)\otimes\Zp=e\,\Zp[\XBq{UU_s}]$.
\item ($B=B_q$.) $\Phicomb(\GDR_q(UU_s))\otimes\Zp\cong\Zp/|X_s|$ carries the action $\varphi\mapsto\varphi\circ U_p^*$ induced from
  $\Hom(X,\Z)$, which is multiplication by $p$; so $e(\Phicomb\otimes\Zp)=0$. This is a non-trivial statement: $v_p|X_s|\ge 2s-2$.
\item (Any $B$.) The Brandt pairing $\langle[x],[y]\rangle=\delta_{xy}$ on $\Zp[X_s]$ restricts to a perfect pairing
  \[e\,\Zp[X_s]\times e^*\Zp[X_s]\to\Zp,\]
  and $\nu$ identifies $e^*\Zp[X_s]$ with $e\,\Zp^{X_s}$ for the \emph{cohomological} $U_p$ ($f\mapsto\sum_jf(\cdot\,\varpi_j)$).
  Thus the perfect pairing is the evaluation pairing $e\,\Zp[X_s]\times e\,\Zp^{X_s}\to\Zp$. $e$ and $e^*$ are different
  idempotents of $\End\Zp[X_s]$ in general.
\end{enumerate}
\stproved{} below ((a)--(c) as in the seed, with the missing equivariance details; (d) with the identification of $e^*$).
\end{lemma}
\begin{proof}
(a) $\deg U_p[x]=\#\{j\}=p$. On $\Zp$ with $U_p=p$, $p^{n!}\to0$. Exactness of $e$ applied to $0\to\Zp[X]^0\to\Zp[X]\to\Zp\to0$
gives $e\Zp[X]^0=e\Zp[X]$. For $H_0$: at $s\ge1$ the level is free (\texttt{if:s02-indef}), $H_0(\Gamma_{UU_s},\Zp)=H_0(|\G|)\otimes\Zp$
for the connected graph $|\G|$, and the isomorphism with $\Zp$ is $(a,b)\mapsto\deg a+\deg b$ on $\coker\partial$; $U_p$ acts
componentwise on $\Zp[X_s]^{2}$ (it commutes with $s_{0*},s_{1*}$ by \texttt{if:s02-quot}(c)), so by $p$ on $H_0$. No density
statement is used: only connectivity of the quotient graph.
(b) $\partial e=v_1-v_0$, so $\partial(U_pe)=\sum_j\partial[e\varpi_j]=p(v_1-v_0)$. Apply $e$ to
$0\to\Xcomb\to\Z[E]\to\Z[V]^0\to0$ and use (a) on $\Z[V]^0\cong\Z$.
(c) With $w=1$, $X=\Z[E]^0$, $\Hom(X,\Z)=\Z^E/\Z\mathbf 1$ and $\Phicomb=\Z^E/(\Z[E]^0+\Z\mathbf1)\xrightarrow{\Sigma}\Z/|E|$. The map
$x\mapsto\langle x,\cdot\rangle$ intertwines $U_p$ with $\varphi\mapsto\varphi\circ U_p^*$, and
$\sum_e\varphi(U_p^*[e])=\sum_{e'}\varphi(e')\cdot\#\{\text{column sum of }U_p\text{ at }e'\}=p\sum_{e'}\varphi(e')$, since
the matrix of $U_p^*$ is the transpose of that of $U_p$ ($w=1$) and every column of $U_p$ has sum $\deg U_p[e']=p$.
$v_p|X_s|$: by \texttt{if:s01-mass}, $|X_s|=\frac{\varphi(q)}{24}\psi(Mp^s)\varphi(Mp^s)$ (for $U=\Uone M$), which has $p$-adic
valuation $\ge2s-2$.
(d) $\langle U_pa,b\rangle=\langle a,U_p^*b\rangle$ gives $\langle e a,b\rangle=\langle a,e^*b\rangle$, so
$e\Zp[X]\perp(1-e^*)\Zp[X]$ and $(1-e)\Zp[X]\perp e^*\Zp[X]$; perfectness of $\delta_{xy}$ on $\Zp[X]$ then gives perfectness on
$e\Zp[X]\times e^*\Zp[X]$. Finally $\nu\circ U_p^*=U_p^{\mathrm{coh}}\circ\nu$ (transpose), whence $\nu e^*=e^{\mathrm{coh}}\nu$.
\end{proof}
\noindent\emph{Why (d) matters.} The family $(e\Zp[X_s])_s$ is an inverse system under $\pi_*$ (stage 01, Lemma Up (b)), and
$(e\Zp^{X_s})_s$ is a direct system under $\pi^*$; the family $(e^*\Zp[X_s])_s$ is compatible with neither $\pi_*$ nor
$\pi^*$ (only with traces, Lemma Up (c)). So the $\Lambda$-adic pairing cannot be the naive limit of (d); it needs the twist of
Def.~\ref{def:s04-BN}.

\begin{proposition}[$e$ kills $H_0$ with $\Vk$-coefficients]\label{prop:s04-ekillsVk}
Let $s\ge1$, $k\ge2$, $q\nmid Dp$, $U$ of level prime to $pq$ with $U_q=\GL_2(\Zq)$. Then $e$ kills $H_0(\Gamma_{UU_s},\Vk)$ (with
$U_p$ induced from $\coker\partial$). Consequently, for every $k$,
\[0\to e\,H_1(\Gamma_{UU_s},\Vk)\to e\,\Brandt{\Uzero{q}UU_s}{\Vk}_*\xrightarrow{\ \partial\ }e\,\Brandt{UU_s}{\Vk}_*^{\,2}\to0\]
is exact, \emph{without} localisation at a non-Eisenstein $\m$.
\end{proposition}
\noindent\sttoprove{} (sketch). By \texttt{if:s02-coh}(c), $H_0=(\Vk)_{\Gamma}$ with $\Gamma=\Gamma_{UU_s}$ dense in
$U^1_p$; this identification is \stpmc{} via strong approximation (\texttt{if:s02-SA}) and must not be replaced by a per-level
check (stage 02, deviation D3). Under it, $U_p$ acts by $v\mapsto\sum_{i=1}^{p}\xi_iv$ for global $\xi_i\in B^\times$ with
$\xi_i\in\varpi_{j(i)}\Kp{s}$ at $p$ (Hecke operator of the commensurator; this transfer of $U_p$ to group homology is the step
still to be written, using \texttt{if:s02-quot}(c)). Every $g=\left(\begin{smallmatrix}a&b\\c&d\end{smallmatrix}\right)\in\varpi_j\Kp{s}$
has $a\equiv c\equiv0$, $d\equiv1\pmod p$, so by the action $P\mapsto P((X,Y)g)$ of stage 01, $gP\equiv P(0,1)\,(bX+dY)^n\pmod p$: modulo $p$, $g$ factors
through $\mathrm{ev}:P\mapsto P(0,1)$ and $\mathrm{ev}(gP)\equiv\mathrm{ev}(P)$. Hence $\mathrm{ev}(U_pv)\equiv p\,\mathrm{ev}(v)\equiv0$ and
$U_p^2\equiv0$ on $(\Vk)_\Gamma/p$, so $U_p$ is topologically nilpotent on the finitely generated $\Zp$-module $(\Vk)_\Gamma$ and $e=0$.
In particular, $e$ kills the \emph{non-zero} Ihara cokernel that stage~03 predicts at level $\Uone{p^s}$ (stage 03, Remark on
T6 (b): $H_0(\GU,\Vk)\ne0$ for every $k$ when $U_p\subseteq\Uone{p^s}_p$): it is non-zero but non-ordinary.
This upgrades the seed's ``with $V_k$ after localisation at a non-Eisenstein $\m$ \stsketch'' (cf.\ stage 03, Remark on T6 (b):
$H_0\ne0$ at $\Uone p$-level, but it is non-ordinary). Note that the analogous argument on $\Brandt{UU_s}{\Vk}$ itself fails (the
$p$ summands live at different points), as it must.

\subsubsection{Hida theory on the definite side}

\begin{definition}[$\Lambda$-adic Brandt module; seed \S7]\label{def:s04-Mord}
For $U$ of level prime to $p$: $\Mord(U):=\varprojlim_{s\ge1}e\,\cOE[X_s(U)]$, limit along $\pi_*$ (which commutes with $U_p$ by
stage 01, Lemma Up (b), and with all $T_\ell,S_\ell,U_\ell,\Pi_\ell,\langle d\rangle$). $\Lam$ acts through the adelic diamonds:
$[u]\cdot a:=\langle u\rangle a$, $u\in1+p\Zp$. The finite group $\mu_{p-1}\times(\Z/M)^\times$ acts as well, and $\Mord(U)$
decomposes into eigenpieces $\Mord(U)^{(\chi)}$.
\end{definition}

\begin{definition}[arithmetic points in the seed's conventions]\label{def:s04-Pk}
For $k\ge2$ put $P^{\rm ad}_k:=\ker\big(\Lam\to\cOE,\ [u]\mapsto u^{-(k-2)}\big)$. The seed's $P_k=\ker([u]\mapsto u^{k-2})$ is the
image of $P^{\rm ad}_k$ under the involution $\iota:[u]\mapsto[u^{-1}]$, i.e.\ it is the right ideal for the $\Lam$-structure via
\emph{classical} diamonds $\langle u\rangle_{\rm cl}=\langle u\rangle^{-1}$.
\end{definition}
\noindent\emph{Justification} \stproved: for $F\in\Brandt{U\Kp1}{\Vk}$ put $f:=\mathrm{ev}\circ F$. Since
$\mathrm{ev}(gP)\equiv d_g^{\,n}\,\mathrm{ev}(P)\pmod{c_g}$ (stage 01, coefficient conventions), $f\bmod p^s$ is a function on $X_s$
with $f(x\delta_d)\equiv d^{-n}f(x)$, i.e.\ the (cohomological, hence also homological) adelic diamond acts by $d^{-(k-2)}$ on the
weight-$k$ piece; and $\mathrm{ev}(\varpi_jP)=\mathrm{ev}(P)$ exactly, so $\mathrm{ev}$ intertwines $U_p$ on $\Vk$-valued forms with
$U_p$ on $\Zp$-valued forms. This matches \lit{hsieh\_2017\_hida-families-and-p-adic-triple-product-l-functions\_arxiv}{\S4.4,
Def.~4.1, p.~31}: there $\langle z\rangle_\Lam x:=x\,\mathrm{diag}(z,z)$, which equals our $\langle z\rangle$ on $X_s$ because
$\mathrm{diag}(z,1)\in\Kp s$, and forms satisfy $f(xz)=f(x)\langle z\rangle^2\langle z\rangle_\Lam^{-1}$. The weight-$k$ point with
trivial character at $p$ also requires $\mu_{p-1}$ to act by $\omega^{-(k-2)}$.

\begin{proposition}[freeness at finite level; seed Prop.~free (a)]\label{prop:s04-freeA}
For $s\ge1$, $e\,\Zp[X_s(U)]$ is a free $\Zp[\Delta_s]$-module. \stproved{} (seed): $\Zp[X_s]$ is $\Zp[\Delta_s]$-free since
$\Delta_s$ acts freely; $U_p$ commutes with $\Delta_s$ (stage 01, Lemma Up (e)), so $e\Zp[X_s]$ is a direct summand, hence
projective, hence free over the local ring $\Zp[\Delta_s]$. Only $\Delta_s$, not $(\Z/p^s)^\times$, is used.
\end{proposition}

\begin{lemma}[Hida's lemma, definite form]\label{lem:s04-hida}
Let $s\ge1$, $\pi':X'_s\to X_s$ the projection ($[\Kp s:\Kpp s]=p$) and $\pi'':X_{s+1}\to X'_s$ (a $\ker(\Delta_{s+1}\to\Delta_s)$-torsor).
\begin{enumerate}[label=(\alph*)]
\item $\Kpp s\varpi_p\Kpp s=\coprod_j\varpi_j\Kpp s$ with the same $\varpi_j$; so $U_p$ is defined compatibly on $X_{s+1},X'_s,X_s$.
\item For $u\in\Kp s$ and each $j$ there are $i$ and $u_j\in\Kpp s$ with $u\varpi_j=\varpi_iu_j$. Hence, cohomologically,
  $U_p\big(\Zp^{X'_s}\big)\subseteq\pi'^*\Zp^{X_s}$, and homologically $U_p(\ker\pi'_*)=0$.
\item $e\pi'_*:e\Zp[X'_s]\to e\Zp[X_s]$ is an isomorphism; and $\pi''_*$ induces $\Zp[X_{s+1}]\otimes_{\Zp[\Delta_{s+1}]}\Zp[\Delta_s]
  \xrightarrow{\sim}\Zp[X'_s]$.
\item Consequently $\pi_*$ induces $e\Zp[X_{s+1}]\otimes_{\Zp[\Delta_{s+1}]}\Zp[\Delta_s]\xrightarrow{\sim}e\Zp[X_s]$, and
  $\rk_{\Zp[\Delta_s]}e\Zp[X_s]$ is independent of $s\ge1$.
\end{enumerate}
\sttoprove{} (all steps elementary; proof below).
\end{lemma}
\begin{proof}[Proof sketch]
(a) as stage 01, Lemma Up (a). (b) For $u=\left(\begin{smallmatrix}a&b\\c&d\end{smallmatrix}\right)\in\Kp s$, the matrix
$\varpi_i^{-1}u\varpi_j$ has lower row $(pc,\,cj+d)$; choose $i\equiv(aj+b)/(cj+d)\bmod p$ to make it integral; then
$pc\in p^{s+1}\Zp$ and $cj+d\equiv1\bmod p^s$, so $u_j\in\Kpp s$. Thus for $f$ invariant under $\Kpp s$,
$(U_pf)(xu)=\sum_jf(x\varpi_iu_j)=(U_pf)(x)$. Dually: if $\pi'_*a=0$ then $\langle U_pa,f\rangle=\langle a,U_pf\rangle=0$ for all
$f$, since $U_pf\in\pi'^*(\cdot)$ is orthogonal to $\ker\pi'_*$. (c) $\pi'_*$ is surjective and $U_p$-equivariant; $e$ is exact and
kills $\ker\pi'_*$. The torsor statement is Lemma stab (b) for $\Delta_{s+1}$ together with $\Kpp s/\Kp{s+1}\cong\ker(\Delta_{s+1}\to\Delta_s)$ via $d$.
(d) Combine (c), using that $e$ commutes with $\otimes_{\Zp[\Delta_{s+1}]}\Zp[\Delta_s]$: $U_p$ commutes with the $\Delta_{s+1}$-action (stage 01, Lemma Up (e)), so $e$ is a $\Zp[\Delta_{s+1}]$-linear idempotent and $e(M\otimes_{\Zp[\Delta_{s+1}]}\Zp[\Delta_s])=(eM)\otimes_{\Zp[\Delta_{s+1}]}\Zp[\Delta_s]$.
\end{proof}
\noindent This is the precise answer to the seed's remark after Prop.~free (``fibres have size $p^2$ while
$\ker(\Delta_{s+1}\to\Delta_s)$ has order $p$''): the extra factor $p$ is the step $X'_s\to X_s$, and $e$ collapses it.

\begin{theorem}[$\Lambda$-freeness and control; seed Prop.~free (b); \texttt{if:s04-Mord}]\label{thm:s04-control}
Let $U$ be of level prime to $p$ and $r:=\rk_{\Zp}e\,\Zp[X_1(U)]$.
\begin{enumerate}[label=(\alph*)]
\item $\Mord(U)$ is a free $\Lam$-module of rank $r$, and for every $s\ge1$,
  $\Mord(U)\otimes_\Lam\cOE[\Delta_s]\xrightarrow{\sim}e\,\cOE[X_s(U)]$, compatibly with all Hecke operators.
\item (weights) For $k\ge2$, $\Mord(U)\otimes_\Lam\Lam/P^{\rm ad}_k$ is dual to $e\,\Brandt{U\Kp1}{\Vk\otimes\cOE}$ (Def.~\ref{def:s04-Pk}).
  On the component where $\mu_{p-1}$ acts by $\omega^{-(k-2)}$ it is dual to $e\,\Brandt{U\cdot\iota_p^{-1}\Uzero{p}_p}{\Vk\otimes\cOE}$, for every $k\ge2$.
  For $k\ge3$ the latter equals $e\,\Brandt{U}{\Vk\otimes\cOE}$ by $p$-stabilisation (\stexp{} in our normalisation).
  \provisional{F4: (b) restated so that the character piece and the level match; before: ``For $k\ge2$ and $\chi$ the component on which $\mu_{p-1}$ acts by $\omega^{-(k-2)}$, $\Mord(U)^{(\chi)}\otimes_\Lam\Lam/P^{\rm ad}_k$ is dual to $e\,M(U\Kp1;\Vk\otimes\cOE)$; for $k\ge3$ one may replace $U\Kp1$ by $U\cdot\iota_p^{-1}\Uzero p_p$ (Hida's control at $\Gamma_0(p)$-level; expected)''; the left side was one $\mu_{p-1}$-eigenpiece and the right side the whole level-$U\Kp1$ module.}
\end{enumerate}
Status: (a) \stpmc{} from Lemma~\ref{lem:s04-hida} and Prop.~\ref{prop:s04-freeA} (choose compatible bases by Nakayama; $\Lam=\varprojlim\cOE[\Delta_s]$);
also \stcited{} \lit{hsieh\_2017\_hida-families-and-p-adic-triple-product-l-functions\_arxiv}{Thm.~4.2 and its proof, p.~32:
$\varprojlim_\alpha e\cO[X_\alpha]$ is $\Lam[\Delta_p]$-free, via [Hid88b, Cor.~10.1]}
(Hsieh's ``Cor.~10.1'' is \emph{Theorem}~10.1 in Hida's paper; there is no Cor.~10.1) and
\lit{longo-vigni\_2011\_control-theorems-for-quaternionic-hida-families\_arxiv}{Prop.~3.6 and its proof, p.~28: $\varprojlim_r
S_2(U_r,\cO_f)^{\rm ord}$ is $\Lam$-free by [Hid88b, Cor.~10.4]}. (b) \stcited{}
\lit{hsieh\_2017\_hida-families-and-p-adic-triple-product-l-functions\_arxiv}{Thm.~4.2(2), p.~32 (forms side; via [Hid88b, Thm.~9.4])};
\lit{longo-vigni\_2011\_control-theorems-for-quaternionic-hida-families\_arxiv}{Thm.~3.7, p.~28 (eigen-components, over a field)};
the dictionary to our $\NUp$, $\Vk$ and $P^{\rm ad}_k$ is \sttoprove{} (Def.~\ref{def:s04-Pk} is the key computation). The
original source is \lit{hida\_1988\_p-adic-hecke-algebras-for-gl2-over-totally-real-fields}{Thm.~8.6, p.~362 (weight independence of
$\mathcal V^{\rm ord}$); Thm.~9.4, p.~369 (control); Thm.~10.1, p.~370 (freeness over $\cO[[Z_a]]$ under $p$-divisibility);
Cor.~10.4, p.~373 ($\Lam$-freeness, $p>2$)}: the theorem headers and page numbers are located in the OCR text, but the statements are
only partly legible, so their exact hypotheses are taken from Hsieh and LV11 (missing/stage-04.md).
\end{theorem}
\noindent\emph{Factor $2$.} At $s\ge1$ there are no stabilisers and $-1\notin UU_s$, so no factor $2$ arises inside the tower.
When comparing with a level containing $-1$ (e.g.\ $U$ itself, $U\cdot\Uzero p_p$, or $\Uzero{q}U$ at $s=0$), the stage 01 rule
applies: $\nu\circ\pi_*=\mathrm{tr}\circ\nu'$ needs $-1\in U'$ or $-1\notin U$, otherwise $\mathrm{tr}\circ\nu'=2\,\nu\circ\pi_*$.
Passing from $U\ni-1$ to $UU_s\not\ni-1$ (e.g.\ from $\hat\cO^\times$ or $\Uzero M$ to $\Uone{Mp^s}$) is exactly the bad case,
and the hypothesis ``$-1\in U'$ or $-1\notin U$'' must be checked at \emph{every} level change used (inside the tower,
$s\ge1$, it holds since $-1\notin UU_{s+1}\subseteq UU_s\not\ni-1$; for $X'_s\to X_s$ likewise); since $2\in\Zp^\times$ this changes nothing over $\Zp$, but every
comparison of pairings across such a step must carry this factor explicitly. Hsieh's $\#\Gamma_{N,x}:=\#(D^\times\cap x\hat R_N^\times x^{-1})\Q^\times/\Q^\times$
equals our $w_x$; the import's \texttt{Nat.card} $=|\Gamma_x|$ is $2w_x$ at levels containing $-1$.

\begin{remark}[GL$_2$ side, for orientation only]\label{rem:s04-hida86}
The seed's pointer ``[Hi86b] Thm.~3.1, Cor.~3.2--3.3'' is to the \emph{elliptic} theory:
\lit{hida\_1986\_iwasawa-modules-attached-to-congruences-of-cusp-forms}{Thm.~3.1, p.~250 (ordinary Hecke algebras free of finite
rank over $\Lam$, $p\ge5$); Cor.~3.2, p.~250 (control, $k\ge2$ for cusp forms); Cor.~3.3, p.~250 (ideals of old forms)}. Hida's
$P_k$ there is normalised via his action (3.1) of $\Zp^\times$ (p.~248) and differs from both $P_k$ and $P^{\rm ad}_k$ by a shift
and/or $\iota$; translate before use. None of these is a statement about $\XB U$.
\end{remark}

\subsubsection{The $\Lambda$-adic definite Ribet sequence}

\begin{theorem}[$\Lambda$-adic definite Ribet sequence; seed Thm.~lambdaribet; \texttt{if:s04-lambdaribet}]\label{thm:s04-lambdaribet}
Let $q\nmid Dp$, $U$ of level prime to $pq$ with $U_q=\GL_2(\Zq)$ ($q$-admissible in the sense of stage 02). Assume $\BG(UU_s,q)$ is connected for every $s\ge1$
(equivalently, Prop.~quot holds at all these levels). This follows from \texttt{Hyp.SA} $\wedge$ \texttt{Hyp.HS}; it is not a finite check, and it does not imply $H_0=V_{U_p^1}$.
\provisional{F7: connectivity hypothesis made explicit; before: no such hypothesis was stated (only the implicit ``$q$-admissible'').}
Put
\[\Xord_{Dq}(U):=\varprojlim_se\,H_1(\Gamma_{UU_s},\Zp)\otimes\cOE\]
(transition maps: corestriction, i.e.\ the restriction of $\pi_*$ on $\Zp[X'_s]$-chains; over $\Zp$, $p\ge5$, $H_1(\Gamma,\Zp)=H_1(\BG,\Z)\otimes\Zp$
by stage 02, D1). Then:
\begin{enumerate}[label=(\alph*)]
\item the sequence of $\Lam$-modules, equivariant for all Hecke operators supported away from $q$ (including $U_p$, $\langle d\rangle$,
  $\Pi_\ell$ for $\ell\mid D$),
  \[0\to\Xord_{Dq}(U)\to\Mord(\Uzero{q}U)\xrightarrow{\ \alpha=(-s_{0*},\,s_{1*})\ }\Mord(U)^{\oplus2}\to0\]
  is exact;
\item all three terms are $\Lam$-free, and $\rk_\Lam\Xord_{Dq}(U)=\rk_\Lam\Mord(\Uzero qU)-2\rk_\Lam\Mord(U)$;
\item at finite level $s\ge1$: $\rk_{\Zp}H_1(\Gamma_{UU_s},\Zp)=(q-1)|X_s(U)|+1$, and $\rk_{\Zp}e\,H_1=\rk e\Zp[X'_{s}]-2\rk e\Zp[X_s]$,
  the ``$+1$'' being the degree line killed by $e$ (here $X'$ means the level $\Uzero qUU_s$).
\end{enumerate}
Status: (a) \stproved{} given \texttt{if:s03-ribet}, Lemma~\ref{lem:s04-ekills}(a) and $\NUp$; (b) \stpmc{} via
Thm.~\ref{thm:s04-control}(a); (c) \stproved{} (seed rank check; \texttt{if:s02-indef}).
\end{theorem}
\begin{proof}
For each $s\ge1$, \texttt{if:s03-ribet}(a) tensored with $\Zp$ gives $0\to H_1\to\Zp[X'_s]^0\to(\Zp[X_s]^0)^2\to0$; $U_p$ commutes
with $s_{0*},s_{1*}$ (\texttt{if:s02-quot}(c), supports $p\ne q$), so $H_1=\ker\alpha$ is $U_p$-stable. Apply the exact functor $e$
and Lemma~\ref{lem:s04-ekills}(a): $0\to eH_1\to e\Zp[X'_s]\to e\Zp[X_s]^2\to0$. The maps commute with $\pi_*$, because
$\NUp$ commutes with $\pi_*$ (stage 01, Lemma Up (b); with the seed's literal $U'_p$ this step fails). All terms are finitely
generated $\Zp$-modules, hence compact, so $\varprojlim^1$ vanishes and $\varprojlim$ is exact (no surjectivity of the transition maps
is needed; Mittag-Leffler for the left-hand system $(eH_1)_s$ holds because its terms are compact). (b): a surjection onto a free module splits, so the kernel is a
direct summand of a free $\Lam$-module, hence free ($\Lam$ local).
\end{proof}
\noindent\emph{Expected rank} \stexp: by \texttt{if:s01-targets} and Thm.~\ref{thm:s04-control}, $\rk_\Lam\Xord_{Dq}(U)$ equals the
$\Lam$-rank of the $q$-new quotient of $\TordB(\Uzero qU)$, i.e.\ of the $Dq$-new ordinary Hecke algebra of $\GL_2$; this is the
rank of the character group of the toric part of Hida's ordinary $Dq$-new family, the shape required by $\HypX$ (stage 05). The
equality ``old part has rank $2\rk\Mord(U)$'' is the rational injectivity of $(s_0^*,s_1^*)$, i.e.\ \texttt{if:s03-ihara}(ii).

\subsubsection{Pairings, Gorenstein property, $\eta^B_F$}

\begin{definition}[Hsieh's twisted $\Lambda$-adic pairing]\label{def:s04-BN}
Let $N=N^+N^-$ with $N^-=D$, $N^+=M$, and $\chi$ a finite-order character mod $Mp$. For $\Lam$-adic forms $f,f'$ (functions on
$X_\infty=B^\times\backslash\Bhat^\times/U_1(Mp^\infty)$, Hsieh Def.~4.1) in $eS^D(N,\chi,\mathbb I)$,
\[B_{N,\alpha}(f,f'):=\sum_{[x]\in X_0(Np^\alpha)}U_p^{-\alpha}f(x\tau^D_{Np^\alpha})\,f'(x)\,\chi_{\mathbb I}(\Nrd x)\,(\#\Gamma_{Np^\alpha,x})^{-1}\pmod{P_\alpha},\qquad
B_N:=\varprojlim_\alpha B_{N,\alpha},\]
with $\tau^D$ the Atkin--Lehner element at $N^+p^\alpha$ and $\chi_{\mathbb I}(z)=\chi(z)\langle\varepsilon_{\rm cyc}(z)\rangle^{-2}\langle\varepsilon_{\rm cyc}(z)\rangle_{\mathbb I}$.
\end{definition}
\noindent\stcited{} \lit{hsieh\_2017\_hida-families-and-p-adic-triple-product-l-functions\_arxiv}{\S4.4, ``A perfect pairing'', and
Def.~4.3, pp.~32--33; finite-level version $\langle f_1,f_2\rangle_N=(f_1,[\tau^D_N]f_2)_N$, \S4.3, p.~30}. $B_N$ is Hecke-equivariant
and perfect; at $k_Q=2$, $B_N(f,f')(Q)=\langle U_p^{-\alpha}f_Q,f'_Q\rangle_{Np^\alpha}$, and in weight $k_Q$ with the extra sign
$(-1)^{k_Q}$ (Hsieh, Lemma~4.4, p.~33). \emph{Translation} \sttoprove: Hsieh's homological $U_p$ on $\cO[X_\alpha]$ is
$x\mapsto\sum_bx\left(\begin{smallmatrix}\varpi_p&b\\0&1\end{smallmatrix}\right)$ (p.~31), i.e.\ $\NUp$, and his $\langle z\rangle_\Lam$ is our
$\langle z\rangle$ (Def.~\ref{def:s04-Pk}); his forms are $\Lam$-duals twisted by $\iota_2$ ((4.6), p.~31). So the seed's ``twist by
$U_p^{-s}w_{p^s}$'' is right in substance, but (i) it lives on a fixed character piece over $\mathbb I$, not on all of $\Mord$, (ii) it
needs the Atkin--Lehner at $M$ as well as at $p^\alpha$, (iii) it carries the norm twist $\chi_{\mathbb I}(\Nrd x)$, and (iv) it
uses $\#\Gamma_{Np^\alpha,x}=\#\big((B^\times\cap x\hat R^\times x^{-1})\Q^\times/\Q^\times\big)$ (Hsieh, (4.1), p.~29), which is our $w_x$
(stage 01: $w_x=|\Gamma_x/\Gamma_x\cap\{\pm1\}|$), i.e.\ $|\Gamma_x|/2$ at $X_0(Np^\alpha)$ where $-1\in\Gamma_x$; the factor $2$ is against
$|\Gamma_x|$ (the import's \texttt{Nat.card}), not against $w_x$. The $\Vk$-pairing of stage 01
(perfect only for $p>k-2$) is \emph{not} used: perfectness comes from the trivial-coefficient tower.

\begin{hypothesis}[Hsieh's residual condition (CR, $\mathrm{supp}\,D$)]\label{hyp:s04-CR}
Let $\m\subset\TordB(U)$, $U=\Uone M$, be a maximal ideal with residual representation $\bar\rho_\m:G_\Q\to\GL_2(\bar\F_p)$
(attached via \texttt{if:s01-targets}). (CR$_D$): $\bar\rho_\m$ is absolutely irreducible (in particular $\m$ is non-Eisenstein),
\emph{$p$-distinguished} ($\bar\rho_\m|_{G_{\Qp}}^{\rm ss}=\chi^+\oplus\chi^-$ with $\chi^+\ne\chi^-$), and ramified at every prime
$\ell\mid D$ with $\ell\equiv1\pmod p$.
\end{hypothesis}
\noindent\stcited{} \lit{hsieh\_2017\_hida-families-and-p-adic-triple-product-l-functions\_arxiv}{\S1.4, Hypothesis (CR,$\Sigma$),
p.~4, with $\Sigma=\mathrm{supp}(N^-)$ as in Thm.~4.5}. Note: $\ell\equiv1$, not $\ell\equiv\pm1$; ``non-Eisenstein'' is implied by
absolute irreducibility and need not be listed separately.

\begin{theorem}[Gorenstein and freeness; \texttt{if:s04-Mord}, second half]\label{thm:s04-Gor}
Assume Hyp.~\ref{hyp:s04-CR} for $\m$, $p>3$, $p\nmid MD$. Then $\Mord(\Uone M)_\m\otimes_\Lam\mathbb I$ (equivalently Hsieh's
$eS^D(N,\mathbb I)_\m$, its $\Lam$-dual) is free of rank one over $\TordB{}_{,\m}$, and $\TordB{}_{,\m}$ (on each character piece) is
Gorenstein. For a primitive family $F$ with $\m_F=\m$, the eigenline $eS^D(N,\mathbb I)[\lambda^D_F]$ is free of rank one over
$\mathbb I$ (its generator $F^D$: the primitive JL lift).
\end{theorem}
\noindent\stcited{} \lit{hsieh\_2017\_hida-families-and-p-adic-triple-product-l-functions\_arxiv}{Thm.~4.5 and its proof, p.~34}.
The proof there cites [Wil95, Thm.~2.1], [Hel07, Cor.~8.11, Rem.~8.12] and [PW11, Prop.~6.4, 6.5]
(\lit{pollack-weston\_2006\_anticyclotomic-mu-invariants-of-modular-forms\_arxiv}{Prop.~6.4, p.~19; Prop.~6.5, p.~20}); Gorenstein
follows from freeness plus the perfect pairing $B_N$. The freeness of $\Mord{}_{,\m}$ is in the \emph{proof}, not the statement, of
Thm.~4.5 (seed cites only ``Thm.~4.5''). Passing between $\Mord$ (homological) and $eS^D$ (forms) uses Hsieh (4.6) and
Thm.~4.2(1); dualising a free rank-one module over a Gorenstein ring is again free of rank one.

\begin{theorem}[$\eta^B_F$ as a self-pairing; \texttt{if:s04-etaB}]\label{thm:s04-etaB}
Under Hyp.~\ref{hyp:s04-CR}, let $F$ be a primitive $\mathbb I$-adic family of tame conductor $N=MD$ with residual ideal $\m$, $F^D$ a
generator of its eigenline (Thm.~\ref{thm:s04-Gor}), $\lambda^D_F:\T^D(N,\chi,\mathbb I)\to\mathbb I$ its eigensystem, and
$\eta^B_F:=\eta(\T^D(N,\chi,\mathbb I)_\m,\lambda^D_F)=\lambda^D_F(\Ann\ker\lambda^D_F)$ (seed \S2.6). Then
\[\eta^B_F=\big(B_N(F^D,F^D)\big)\subseteq\mathbb I .\]
Moreover $(\eta^B_F)\supseteq(\eta_F)$ (the $\GL_2$ congruence ideal), with equality if $\bar\rho_\m$ is ramified at all $\ell\mid D$.
\end{theorem}
\noindent\stcited{} \lit{hsieh\_2017\_hida-families-and-p-adic-triple-product-l-functions\_arxiv}{Def.~4.12, p.~43 ($\eta_{F^D}:=B_{N_F}(F^D,F^D)$);
Remark~7.8, p.~71 (it generates the congruence ideal of $\lambda_f:\T^D(N,\psi,\mathbb I)\to\mathbb I$ under (CR,$\Sigma^-$), and
$(\eta_{f^D})\supseteq(\eta_f)$, equality if ramified at all $\ell\in\Sigma^-$)}. The general principle (free rank one over Gorenstein
+ perfect equivariant pairing $\Rightarrow$ congruence ideal $=$ self-pairing of a generator) is \sttoprove{} as pure commutative
algebra (Lean-friendly; cf.\ \lit{hsieh\_2017\_hida-families-and-p-adic-triple-product-l-functions\_arxiv}{Remark~7.3, p.~67: $T$ Gorenstein,
$H$ free of rank $d$ with perfect equivariant pairing $\Rightarrow C(\lambda)^d=(\det\langle e_i,e_j\rangle)$}). This is a statement about $B$
alone; the comparison with $\eta_F$ is the input to stage 06.

\subsubsection{Nearly ordinary versus ordinary (seed \S3.3)}
\begin{remark}\label{rem:s04-nearly}
On the definite side the nearly ordinary tower uses $U^{(p)}\times\{g\in\GL_2(\Zp):c\equiv0\bmod p^s\}\cap\ker(\text{torus}\to(\Z/p^s)^{\times2})$
and $\Lam_2=\cOE[[(1+p\Zp)^2]]$. Twisting by $x\mapsto\psi(\tau(x))$ ($\tau$ the norm character of stage 03, (C5); $\psi$ a character of
$1+p\Zp$) identifies the nearly ordinary module with $\Mord\hat\otimes_{\cOE}\cOE[[1+p\Zp]]$, the second factor acting by the central
(Nrd-)character. \stsketch{} (the seed's claim, made explicit; it uses $\Q_{>0}\backslash\hat\Q^\times/\Zhat^\times=1$). Nothing in
stages 04--06 requires the nearly ordinary theory.
\end{remark}

%----------------------------------------------------------------------
\subsection{Proof strategy}
%----------------------------------------------------------------------
\begin{enumerate}
\item \emph{$e$ as an algebraic object.} Def.~\ref{def:s04-e} for finite and finitely generated modules; exactness is formal from
  $e$ being an idempotent in a commutative algebra acting on all terms.
\item \emph{Lemma ekills.} Pure degree counting with $\NUp$ (each $U_p[x]$ has $p$ terms) and stage 03's graph sequences; (d) is
  linear algebra of adjoint idempotents. No density or strong approximation is used for $V=\Zp$.
\item \emph{Control.} Lemma~\ref{lem:s04-hida}: the matrix identity $u\varpi_j=\varpi_iu_j$ ($u_j\in\Kpp s$), plus the free
  $\Delta$-torsor $X_{s+1}\to X'_s$. Then Nakayama for the limit. This replaces ``Hida's argument (cited)'' by a self-contained proof
  for trivial coefficients; the weight-$k$ comparison (Thm.~\ref{thm:s04-control}(b)) stays cited, with the computation of
  Def.~\ref{def:s04-Pk} as the dictionary.
\item \emph{Thm.~lambdaribet.} Exactness of $e$ and $\varprojlim$; freeness from control.
\item \emph{Hsieh inputs.} Cited; our work is the dictionary (conventions match on $U_p$ and $\Lam$; differ on $\#\Gamma$ versus $w_x$,
  on forms versus homology, and on $U_q$ for $q\mid D$: Hsieh's sign for $U_q$, $q\mid D$, is not imported, cf.\ stage 01, Example $\Pi_{11}$).
  \provisional{F6/D3: replaced the clause ``where $\mathrm{JL}(U_q)=-U_q$''; before: ``and on $U_q$ for $q\mid D$, where $\mathrm{JL}(U_q)=-U_q$, cf.\ stage 01, Example $\Pi_{11}$''.}
\item \emph{No use of $\Psi^*\Psi$.} This stage does not use the $\GL_2$ amalgam. Should a $\Lam$-adic level-raising module be
  needed later, only $\det(\Psi^*\Psi)=(q+1)^2-q^{2-k}T_q^2$ and $\coker(\Psi^*\Psi)_\m$ may be used (stage 03, D4: the
  off-diagonal normalisation is not canonical).
\item \emph{Prop.~\ref{prop:s04-ekillsVk}.} Needs the transfer of $U_p$ to the group homology of $\GU$; low priority (not needed for
  the $\Lam$-adic statements, which use trivial coefficients only).
\end{enumerate}

%----------------------------------------------------------------------
\subsection{Dependencies and interfaces}
%----------------------------------------------------------------------
\emph{Imported} (restated verbatim from SPLIT.md):

\begin{interface}[\texttt{if:s01-classset}]
$X_B(U)$, $\Gamma_x$ and $w_x$ as defined in the conventions above. Facts: $X_B(U)$ is finite, $\Gamma_x$ is finite,
$\Nrd(U)=\Zhat^\times$ for all the standard levels, and Lemma stab (a)--(c) holds. \status{imported from stage 01}
\end{interface}
\begin{interface}[\texttt{if:s01-brandt}]
$M(U;V)$ together with the Brandt pairing $\langle[x],[y]\rangle=w_x\delta_{xy}$, the norm-map identification
$\Zp[X]\cong\Zp^X$ when $p\nmid|\Gamma_x|$, and the right-translation maps $x\mapsto xg$ for $g\in B_q^\times$.
\status{imported from stage 01}
\end{interface}
\begin{interface}[\texttt{if:s01-hecke}]
The Hecke operators $T_\ell$, $S_\ell$, $U_\ell$, $\langle d\rangle$, $U_p$ (with fixed normalisation $N_{U_p}$) and $\Pi_\ell$
($\ell\mid D$). All operators supported away from $q$ commute with right translation by $B_q^\times$. $T^*=T_{g^{-1}}$; $T_\ell$ is self-adjoint only at level $\Uzero M$.
\provisional{F15: exporter's sentence from stage 01; before: ``Self-adjointness for the pairing''.}
\status{imported from stage 01}
\end{interface}
\begin{interface}[\texttt{if:s01-mass}]
The mass formula $\sum_x1/|\Gamma_x|=\frac1{24}\prod_{\ell\mid D}(\ell-1)$ at maximal level (the full-unit convention, as in the import's
\texttt{Hyp.EichlerMass}), and its level change. \status{imported from stage 01}
\end{interface}
\begin{interface}[\texttt{if:s01-targets}]
Thm JLclassical (cited) and the $\Lambda$-adic Hecke JL $\T^{\ord,D\text{-new}}\cong\T^{\ord}_B$ ([LV], [Hs]; cited).
\status{imported from stage 01}
\end{interface}
\begin{interface}[\texttt{if:s02-SA}]
Strong approximation: $\hat B^1=B^1(B^1_q\times U^{1,(q)})$ for $q\nmid D$, with $U$ as in Prop quot. Status: cited (Voight Ch.~28), or a
hypothesis in Lean. \status{imported from stage 02}
\end{interface}
\begin{interface}[\texttt{if:s02-quot}, \texttt{if:s02-coh}, \texttt{if:s02-indef}]
Prop quot (a)--(c): the bijections, the degeneracies $s_0[x]=[x]$ and $s_1[x]=[x\alpha]$, stabilisers, and Hecke equivariance.
Prop coh (the two-term complex) and Prop harm. $H_0(\Gamma_U,V)=V_{U_p^1}$. The definition $\cX^{(q)}_{Dq}(U;V):=H_1(\Gamma_U,V)$.
The rank at a free level is $(q-1)|X_B(U)|+1$. Over $\Z$, $\cX^{(q)}_{Dq}(U)=H_1(\BG(U,q),\Z)$. $H_0=V_{U_p^1}$ needs full strong approximation. \status{imported from stage 02}
\provisional{F14: the two amendments of the exporter added to the import block; before: block ended after ``The rank at a free level is $(q-1)|X_B(U)|+1$'', and $H_0(\Gamma_U,V)=V_{U_p^1}$ was stated without the strong-approximation caveat.}
\end{interface}
\begin{interface}[\texttt{if:s03-graph}, \texttt{if:s03-ribet}, \texttt{if:s03-ihara}]
Weighted graphs, $X(\G)$, $\langle,\rangle_w$, $\Phicomb$, Prop phi; the DR graph $\GDR_q(U)$ and $\Xcomb_q(U)$. Thm ribet (a) over
$\Z$ and (b) with coefficients, with $\alpha=(-s_{0*},s_{1*})$ (exact sign convention). Cor ihara (i)--(iii). (iii) says $H_0$ is
Eisenstein; its status is sketched. \status{imported from stage 03}
\end{interface}
\noindent How they are used, with the stage 01/02 precisions: \texttt{if:s01-mass} (level change) only for $v_p|X_s|$ in
Lemma~\ref{lem:s04-ekills}(c); \texttt{if:s02-SA} only in Prop.~\ref{prop:s04-ekillsVk}; $\NUp$ and Lemma Up (b) in Def.~\ref{def:s04-Mord} and
Thm.~\ref{thm:s04-lambdaribet}; ``self-adjointness'' only in the corrected form $U_p^*=\langle p\rangle_MU'_p$ (Lemma~\ref{lem:s04-ekills}(d));
$\nu$ is the identity at $s\ge1$; Lemma stab (b) for $\Delta_s$ only; $H_0=V_{U^1_p}$ only via \texttt{if:s02-SA} (Prop.~\ref{prop:s04-ekillsVk}),
never by a per-level check; over $\Z$, $\cX^{(q)}=H_1(\BG(U,q),\Z)$ (stage 02, D1), which agrees with $H_1(\GU,\cdot)$ over $\Zp$.

\emph{Exported:}
\begin{interface}[\texttt{if:s04-ekills}]\label{if:s04-ekills}
Lemma~\ref{lem:s04-ekills} (a)--(c), plus (d): perfect pairing $e\Zp[X_s]\times e^*\Zp[X_s]\to\Zp$, equivalently the evaluation pairing
$e\Zp[X_s]\times e\Zp^{X_s}$ (cohomological $e$). Status \stproved. Addendum: Prop.~\ref{prop:s04-ekillsVk} ($e$ kills $H_0$ with $\Vk$), \sttoprove.
\end{interface}
\begin{interface}[\texttt{if:s04-Mord}]\label{if:s04-Mord}
$\Mord(U):=\varprojlim_se\,\Zp[X_B(UU_s)]$ ($UU_s=U^{(p)}\times\Uone{p^s}_p$, pushforward, $U_p=\NUp$), with $\Lam$ acting by adelic
diamonds and arithmetic ideals $P^{\rm ad}_k=\ker([u]\mapsto u^{2-k})$ (Def.~\ref{def:s04-Pk}). $\Lam$-freeness and control
(Thm.~\ref{thm:s04-control}; \stpmc{} by Lemma~\ref{lem:s04-hida}; cited [Hs, Thm.~4.2], [LV11, Prop.~3.6, Thm.~3.7], Hida 1988).
Under Hyp.~\ref{hyp:s04-CR} (CR$_D$: $\bar\rho_\m$ absolutely irreducible, $p$-distinguished, ramified at all $\ell\mid D$ with
$\ell\equiv1\bmod p$), $\TordB{}_{,\m}$ is Gorenstein and $\Mord{}_{,\m}$ is free of rank one over it (Thm.~\ref{thm:s04-Gor}; cited [Hs, Thm.~4.5 and proof]).
\end{interface}
\begin{interface}[\texttt{if:s04-lambdaribet}]\label{if:s04-lambdaribet}
Thm.~\ref{thm:s04-lambdaribet}; $\Xord_{Dq}(U):=\varprojlim_se\,H_1(\Gamma_{UU_s},\Zp)$, $\Lam$-free of rank
$\rk\Mord(\Uzero qU)-2\rk\Mord(U)$. Status: exactness \stproved, freeness \stpmc. Hypothesis: $\BG(UU_s,q)$ is connected for every $s\ge1$ (equivalently Prop.~quot at all these levels); it follows from \texttt{Hyp.SA} $\wedge$ \texttt{Hyp.HS}, is not a finite check, and does not imply $H_0=V_{U_p^1}$.
\provisional{F7: hypothesis added to the interface; before: no connectivity hypothesis stated.}
\end{interface}
\begin{interface}[\texttt{if:s04-etaB}]\label{if:s04-etaB}
Thm.~\ref{thm:s04-etaB}: $\eta^B_F=(B_N(F^D,F^D))$ for Hsieh's twisted pairing $B_N$ (Def.~\ref{def:s04-BN}) on the $\chi$-piece over
$\mathbb I$, under Hyp.~\ref{hyp:s04-CR}; and $(\eta^B_F)\supseteq(\eta_F)$ with equality if $\bar\rho_\m$ is ramified at all $\ell\mid D$.
Status \stcited{} [Hs, Def.~4.12, Rem.~7.8].
\end{interface}

\paragraph{Interface deviations.}
\begin{enumerate}[label=(D\arabic*)]
\item \texttt{if:s04-Mord}: SPLIT/seed write $P_k=\ker([u]\mapsto u^{k-2})$ with $\Lam$ acting by (adelic) diamonds. In the stage 01
  conventions the weight-$k$ point is $P^{\rm ad}_k=\ker([u]\mapsto u^{2-k})$; equivalently keep the seed's $P_k$ but let $\Lam$ act by
  classical diamonds $\langle u\rangle^{-1}$. We export $P^{\rm ad}_k$ and the adelic action (which is Hsieh's).
\item \texttt{if:s04-Mord}: ``control (cited)'' is upgraded to \stpmc{} by an elementary proof (Lemma~\ref{lem:s04-hida}) for trivial
  coefficients. The meaning of $UU_s$ is fixed ($U$ of level prime to $p$). Hsieh's residual condition is stated explicitly
  (Hyp.~\ref{hyp:s04-CR}); it already implies non-Eisenstein.
\item \texttt{if:s04-ekills}: (d) is added in the precise form ($e^*$ versus $e$; it is the cohomological $e$ under $\nu$); and the
  $\Vk$-addendum Prop.~\ref{prop:s04-ekillsVk}.
\item \texttt{if:s04-etaB}: the pairing is Hsieh's $B_N$ on a fixed-character piece over $\mathbb I$, with Atkin--Lehner at $Mp^\alpha$,
  norm twist and $\#\Gamma$ weights ($\#\Gamma_{Np^\alpha,x}=w_x=|\Gamma_x|/2$ at levels containing $-1$; revised after the stage-04 audit); and the comparison $(\eta^B_F)\supseteq(\eta_F)$ is
  added for stage 06.
\item Factor $2$: wherever a comparison crosses from a level containing $-1$ to $UU_s$, $\mathrm{tr}\circ\nu'=2\,\nu\circ\pi_*$ (stage 01);
  harmless over $\Zp$ but recorded.
\end{enumerate}

%----------------------------------------------------------------------
\subsection{Lean, HACS, talk}
%----------------------------------------------------------------------
\paragraph{Math track (independent).}
M1: write the proof of Def.~\ref{def:s04-e} (finite monoid idempotent; exactness). M2: write Lemma~\ref{lem:s04-hida} in full
(matrix identity, torsor, Nakayama limit) and Thm.~\ref{thm:s04-control}(a). M3: the dictionary to Hsieh (forms versus homology via
(4.6), $\iota_2$, $\#\Gamma$ versus $w_x$, Atkin--Lehner) and to LV11; settle the weight-$k$ statement (b) in our conventions.
M4: the commutative-algebra lemma ``free rank one over Gorenstein + perfect pairing $\Rightarrow$ congruence ideal $=$ self-pairing''.
M5 (low priority): Prop.~\ref{prop:s04-ekillsVk} (transfer of $U_p$ to $H_*(\GU,\cdot)$). M6: Remark~\ref{rem:s04-nearly} in detail.

\paragraph{Lean track (independent).}
\begin{itemize}
\item L1 (easy): ordinary idempotent for an endomorphism of a finite module (\texttt{Module.End}, finite monoid, $x^{n!}$ stabilises);
  exactness of $e$ on short exact sequences of $\Z/p^m[T]$-modules; topologically nilpotent $\Rightarrow e=0$.
\item L2 (easy): Lemma~\ref{lem:s04-ekills}(a)--(d) for an abstract finite set $X$ with a ``$U_p$'' given by a $\Z_{\ge0}$-matrix of
  constant column sum $p$; (b), (c) on top of stage 03's weighted-graph files (incidence matrices).
\item L3 (medium): Lemma~\ref{lem:s04-hida}(b) as a $2\times2$ matrix statement over $\Zp$; (c), (d) for abstract free torsors.
\item Hypotheses (\texttt{Prop}s, never axioms): \texttt{Hyp.HidaControl} (Thm.~\ref{thm:s04-control}(b)), \texttt{Hyp.HsiehGor}
  (Thm.~\ref{thm:s04-Gor}), \texttt{Hyp.HsiehEta} (Thm.~\ref{thm:s04-etaB}); (a) of Thm.~\ref{thm:s04-control} can be proved, not assumed,
  once L3 exists. Thm.~\ref{thm:s04-lambdaribet}(a) at each finite level from stage 03's Lean statement plus L1.
\item The $\Lam$-adic Thm.~\ref{thm:s04-lambdaribet} takes \texttt{$\forall s\ge1$, BGConnected (U U$_s$) q} or \texttt{Hyp.SA $\wedge$ Hyp.HS} as an argument; both remain registered hypotheses (human decision D5), since no finite check discharges them.
  \provisional{F7/D5: Lean-track line added; before: no such line.}
\end{itemize}

\paragraph{HACS track (independent).}
\begin{itemize}
\item T3: $B_3$, tree at $2$, $p=5$, $s=1,2,3$ ($|S_s|=2,50,1250$), mod $5^3$: compute $U_p=\NUp$ on $\Z/125[S_s]$, $e=U_p^{N}$, check
  Lemma~\ref{lem:s04-ekills}(a), (d), and Lemma~\ref{lem:s04-hida}(d) ($\rk_{\Z/125[\Delta_s]}e$ constant in $s$), and the
  $\Lam$-adic Ribet ranks (Thm.~\ref{thm:s04-lambdaribet}(b),(c)).
\item H6: compare with the seed's literal $U'_p$: Lemma~\ref{lem:s04-hida}(d) must \emph{fail} for $U'_p$ with $\pi_*$ (sanity check of $\NUp$).
\item C2 (seed): ordinary ranks $\rk e\Zp[X_1]$ versus the number of ordinary $D$-new forms of weight $2$, level $\Gamma_1(Mp)\cap\Gamma_0(D)$
  (expected equal; Eisenstein line killed); and weight-$k$ control via Def.~\ref{def:s04-Pk} ($\langle d\rangle=d^{2-k}$).
\item H7: for $D=q$, $\Phicomb\otimes\Z_5=\Z/5^{v_5|S_s|}$ and $e$ acting as $0$ on it (Lemma~\ref{lem:s04-ekills}(c)).
\end{itemize}

\paragraph{Talk track (independent).}
Slide ``integral subtleties are non-ordinary'': Mazur's $\mathrm{num}\frac{q-1}{12}$ and the Eisenstein cokernels have $U_p=p$ (each
point has $p$ neighbours under $U_p$), so $e=\lim U_p^{n!}$ kills them; what survives is a free $\Lam$-module and an exact
$\Lam$-adic Ribet sequence. One picture: the tower $S_{s+1}\to X'_s\to S_s$ with the factor $p$ absorbed by $e$.

%----------------------------------------------------------------------
\subsection{Risks}
%----------------------------------------------------------------------
\begin{enumerate}[label=R\arabic*.]
\item \emph{Normalisations} (highest). $P_k$ versus $P^{\rm ad}_k$, adelic versus classical diamonds, $\NUp$ versus Hsieh/LV/Hida;
  a silent mismatch flips weights $k\leftrightarrow 4-k$ in specialisations. Mitigation: Def.~\ref{def:s04-Pk} and HACS C2.
\item \emph{Hsieh's freeness} rests on [Wil95], [Hel07] and [PW11] (Helm's paper is in literature/ only as a 2004 preprint, not
  text-verified here); the residual condition (CR$_D$) excludes $\ell\mid D$, $\ell\equiv1\bmod p$, $\bar\rho$ unramified.
  \provisional{F20: deleted ``which is exactly where stage 06's defect is non-trivial'' ((CR$_D$) excludes only $\ell\equiv1$ with $\bar\rho$ unramified, which (H0) excludes anyway; the defect is non-trivial whenever $\bar\rho$ is unramified at $q$).}
  Stage 06 must check that its standing hypotheses are compatible with
  Hyp.~\ref{hyp:s04-CR}.
\item \emph{Pairing twist.} The $\Lam$-adic pairing is only cited; its definition in our conventions (Atkin--Lehner elements, norm twist,
  factor $2$, sign $(-1)^k$) is \sttoprove. Stage 05 (T7, Ohta twist) depends on it.
\item \emph{Prop.~\ref{prop:s04-ekillsVk}} depends on transferring $U_p$ to group homology of $\GU$ and on strong approximation
  (density); keep it off the critical path.
\item \emph{Hida 1988 only partly legible} (OCR): theorem numbers and pages located, hypotheses taken from Hsieh and LV11.
\end{enumerate}
